\chapter{CAD-Free Model-Based Blending with Residual RL}
\label{chap:blending}

\section{Blending Objective}
The blending stage converts the residual-removal map into structured additional processing. Its objective is to reduce residual error in the core, follow the decaying transition target, preserve intact surroundings, and maintain stable contact. It is a model-based restoration skill: the spatial target and process demand have already been defined by Chapter~\ref{chap:removal_estimation}.

\section{Local Coverage-Path Generation}
The planner generates a local two-dimensional raster or contour path over cells where \(\Delta h(q)\) exceeds a declared threshold. It does not require a complete CAD model or full 3D reconstruction. The local parameterization provides coverage direction and spatial indexing; unknown height is handled by admittance, and moderate orientation variation is handled by residual RL.

Path topology, boundary offset, turning strategy, and collision/reachability checks are \todoexp{TBD}. The planned scope is limited to local, smooth, continuously varying surfaces without sharp edges, steps, grooves, or discontinuities.

\section[Nominal Force, Velocity, Overlap, and Pass Planning]{Nominal Force, Velocity, Overlap,\\ and Pass Planning}
Using the calibrated process model, the planner assigns nominal desired contact force \(f_{\mathrm{nom}}^{\mathrm{des}}\), tangential feed velocity \(v_{\mathrm{nom}}(q)\), path spacing, overlap, and pass count. These quantities should be chosen together because they jointly determine dwell and accumulated removal. A conservative optimizer or rule-based inversion of Eq.~\eqref{eq:preston_spatial} may be used; the implemented method is \todoexp{TBD}.

\section{Admittance-Based Surface-Height Adaptation}
The planner produces nominal Cartesian position \(p_{\mathrm{nom}}\). Admittance control changes its normal component:
\begin{equation}
  p_{\mathrm{cmd}}=p_{\mathrm{nom}}+\Delta p_{\mathrm{adm}},
  \label{eq:blending_admittance}
\end{equation}
so that unknown surface-height variation is accommodated while tracking the desired contact force. The desired contact force remains a controller-level reference. Tangential coverage direction is not generated by admittance.

\section{Residual Tool-Orientation Adaptation}
During blending only, bounded roll and pitch corrections adapt the tool to smooth unseen curvature:
\begin{equation}
  R_{\mathrm{cmd}}
  =
  R_{\mathrm{nom}}
  \operatorname{Exp}\!\left(
    [\Delta\phi_t,\Delta\theta_t,0]_\times
  \right).
  \label{eq:orientation_residual}
\end{equation}
Yaw and coverage direction remain the planner's responsibility unless the final implementation establishes otherwise. The policy must not assume online access to a CAD-derived or perfect surface normal. A simulated normal may be used as privileged reward information during training only if that separation is stated.

\section[Residual-Removal-Aware Feed-Velocity Adaptation]{Residual-Removal-Aware Feed-Velocity\\ Adaptation}
The velocity command is
\begin{equation}
  v_{\mathrm{cmd}}
  =
  \operatorname{clip}
  \left(v_{\mathrm{nom}}(q_t)+\Delta v_t,
  v_{\min},v_{\max}\right).
  \label{eq:velocity_residual}
\end{equation}
The objective is to achieve spatially varying target removal at the desired locations. The policy should slow where residual demand is large, move faster where little removal remains, reduce unnecessary processing near target completion, and slow when force, contact, or orientation becomes unstable. Feed velocity is not called true sliding velocity; the latter is reserved for actual tool--surface relative motion in the removal model.

\section{Residual Policy Observation, Action, and Reward}
The residual action is
\begin{equation}
  \Delta a_t^{\mathrm{RL}}
  =
  [\Delta\phi_t,\Delta\theta_t,\Delta v_t]^\mathsf{T}.
  \label{eq:residual_action}
\end{equation}
A conceptual observation may include local residual-removal demand, six-axis wrench, force-tracking error, current orientation, current and nominal feed velocity, and local path state:
\begin{equation}
  o_t^{\mathrm{RL}}
  =
  [\Delta h_{\mathrm{local}},F_x,F_y,F_z,\tau_x,\tau_y,\tau_z,
  e_F,R_t,v_t,v_{\mathrm{nom}},\xi_t].
  \label{eq:residual_observation}
\end{equation}
This is not an implemented state specification. \todoexp{Finalize the RL observation, normalization, state estimator, algorithm, network, and update rate.}

A conceptual reward may penalize residual-removal error, force-tracking error, tangential wrench, action magnitude, action change, contact loss, and safety violation:
\begin{equation}
  r_t =
  -c_h e_h
  -c_F e_F^2
  -c_w\|w_t^{\mathrm{tan}}\|^2
  -c_a\|\Delta a_t^{\mathrm{RL}}\|^2
  -c_s
  -c_{\mathrm{safety}}.
  \label{eq:conceptual_rl_reward}
\end{equation}
Terms, weights, temporal credit assignment, privileged signals, and termination are \todoexp{TBD}; Eq.~\eqref{eq:conceptual_rl_reward} is only a design envelope.

\begin{figure}[t]
  \centering
  \fbox{\begin{minipage}[c][0.23\textheight][c]{0.90\textwidth}
    \centering
    local 2D nominal coverage path\\
    \(+\ \Delta p_{\mathrm{adm}}\) for surface height\\
    \(+\ [\Delta\phi,\Delta\theta]\) for smooth orientation variation\\
    \(+\ \Delta v\) for residual-removal-aware tangential feed velocity\\
    \(\downarrow\) final bounded Cartesian command
  \end{minipage}}
  \caption{Blending command composition. The residual policies modify model-based blending and never correct the demonstrated IL trajectory.}
  \label{fig:blending_command}
\end{figure}

\section{Interaction between Orientation and Velocity Policies}
Orientation changes contact distribution and measured wrench, while velocity changes dwell, force transients, and removal accumulation. Their coupled effects require a common observation convention, compatible update rates, and final arbitration by safety limits. Experiment~V compares nominal execution, classical moment-based orientation adaptation, orientation-only residual RL, and the full orientation--velocity residual. Separate and joint policies remain \todoexp{to be determined from the final implementation}.

\section{Iterative Removal-Map Update and Replanning}
After blending pass \(k\), the estimated increment can update the accumulated map:
\begin{align}
  \hat h_{k+1}(q)
  &=\hat h_k(q)+\Delta\hat h_{\mathrm{blend},k}(q),\\
  \Delta h_{k+1}(q)
  &=\max\{0,h^*(q)-\hat h_{k+1}(q)\}.
  \label{eq:iterative_update}
\end{align}
This defines the loop \emph{Plan \(\rightarrow\) Execute \(\rightarrow\) Update Removal Map \(\rightarrow\) Replan}. \todoexp{Iterative replanning is a planned extension unless verified in the final implementation}; it must not be reported as completed before that verification.

\section[Safety Limits and Expected Failure Modes]{Safety Limits and Expected\\ Failure Modes}
Actions are clipped before command composition. Deterministic limits cover normal and tangential force, moments, feed velocity, orientation correction, correction rate, admittance displacement, workspace, contact loss, cumulative predicted removal, and emergency stop. Expected failures include TIF mismatch, incorrect residual demand, policy saturation, oscillatory orientation or velocity, contact loss, conflict between residual channels, and extrapolation beyond trained curvature. Interventions and saturation frequency are reported as outcomes rather than silently discarded.
